\section{Discussion}

We assessed methods for segmenting diabetic retinopathy lesions in fundus photographs by auditing published numbers, scoring and rerunning the original models, retraining \rn{n.main} models under one protocol, and testing all of them on sources they had never seen. The design separates three things that a single leaderboard number mixes: what a paper reports, what its code produces, and what its architecture contributes once data, budget and evaluation are equal.

\paragraph{Reproducibility (Q1).} Most fundus-specific methods cannot be rerun: \rn{funnel.train} of \rn{funnel.dr} have training code. Where code and weights exist, the authors' evaluation code reproduced the released numbers exactly, so the distance between a paper and our protocol is a property of the metric and not of the model; but the released training recipe, run with the authors' settings on three seeds, produced models 5--7 points below the released weights. The literature audit shows why published tables cannot be read as rankings: most papers do not say how the AUPR was computed, the fixed-threshold implementation adds a model-dependent amount, and nearly half of the baseline numbers (\rn{lit.base.copied1}\% identical to the method's own paper, \rn{lit.base.copied2}\% to an earlier copy) were carried from one paper to the next, at times with transposed columns.

\paragraph{Comparison (Q2).} Under equal conditions two adapted foundation models (a low-rank adaptation of SAM and RETFound with a segmentation head) and the general-purpose HRNet-W48 were at the top of both internal tests, SAM-Adapter, Swin-Unet and the self-configuring nnU-Net followed within 0.04, and the seven fundus-specific networks were below them, in line with what has been found for three-dimensional medical segmentation\cite{isensee2024revisited}; over the eight test sets \rn{cd.ntop} of the \rn{n.ranked} ranked models could not be separated from the first. That statement holds for the budget and the 1,440-pixel canvas of the unified protocol: the strongest fundus-specific method, HRDecoder, matched the leading models when run with its authors' released weights, which were trained at twice the resolution and on several times the pixels, and a lesion-specific decoder may need exactly the conditions a shared protocol takes away. The differences among the leading models were of the size of the seed-to-seed variation, whereas the loss of a single rare class, or a recipe that does not converge within the budget, moved a model by ten times as much. Outside the training source the order was less stable than the internal tables suggest: internal rank predicted independent rank on the sources whose lesions are drawn the same way, and the lead on coarsely labelled, region-labelled or ultra-widefield images went to models from the lower half of the internal tables. A model chosen on IDRiD or DDR alone is therefore chosen for the labelling style of that source as much as for the lesions. Where one score is needed, averaging the probability maps of several retrained models gained 0.02--0.06 \mAUPR{} over the best single model on every fine-grained test set and $-$0.01 to 0.02 on the coarse and region-labelled ones, more than the differences among the leading architectures; the four adapted foundation models averaged by themselves gave the same gain. On the one source with several independent expert annotations, the agreement of the better models with the expert consensus was within the agreement among the experts, with the caveats that a region consensus is a coarse ground truth and that the experts' confidence maps give only a four-point curve.
The strongest single models of the study were not in the ranked table. Under one decoder head, one recipe and one training amount, RETFound's two encoders scored no higher than the natural-image encoders they were initialised from ($-$\rn{fm.enc.c1.int.abs} and $-$\rn{fm.enc.c2.int.abs} \mAUPR{} on DDR, $-$\rn{fm.enc.i.c1.int.abs} and $-$\rn{fm.enc.i.c2.int.abs} on IDRiD, one seed each), and the ImageNet masked autoencoder and DINOv2, ViT-L encoders that have never seen a fundus photograph, reached \rn{fm.enc.mi.int} and \rn{fm.enc.di.int} on DDR and \rn{fm.enc.i.mi.int} and \rn{fm.enc.i.di.int} on IDRiD under that head, above every ranked model on both sources. These cells were trained as controls after the ranked rows had been registered, each with one seed and one learning rate, so they do not enter the ranks; but they say that the lead of a ViT-L encoder with a UPerNet head over the convolutional and fundus-specific networks is a lead of the architecture and of the scale of pretraining behind the encoder, which retinal pretraining did not add to, and that the leaderboard of this task would be headed by a generic encoder and a generic decoder if such a configuration were registered. Confirming this with more seeds and a learning-rate sweep is the first item of future work. Zero-shot use of promptable models is not a baseline for this task, and results obtained with prompts from the ground truth should be labelled as upper bounds.

\paragraph{What decides performance (Q3).} The largest effects came from the data and the protocol rather than from the architecture. Whether the test images share the annotation style and the field of view of the training images moved \mAUPR{} by 0.16 on average and up to 0.28 between the four training sources, as much as the differences between the architectures of the main table: models trained on DDR kept or improved their score on the other pixel-labelled standard-field sources and lost a third to a half of it where the labels are coarse or regional or the field is ultra-wide, and models trained on the 44 IDRiD images lost on every other source, as Playout and Cheriet observed for U-Net-type baselines\cite{playout2024crossdatasetgeneralization}. Mixed or coarse labels added to fine ones did not help and often hurt, and the most useful foreign source was the largest fine-grained one rather than the largest one. Shrinking the canvas from 1,440 to 512 pixels cost up to 0.07 \mAUPR{} and between a third and three fifths of the small lesions, more images of the same source helped with diminishing returns, and merging the fine-grained sources improved the scores outside them. An internal test therefore says little about new images unless they are labelled like the training set, a conclusion shared with assessments in other fields\cite{lu2026tcrepitopenaturemethods}.

\paragraph{Recommendations.} For developers of methods: (1) use the official split, and publish the list of image identifiers when a new split is needed; (2) report the exact pooled AUPR at the original resolution inside the field of view, and name the implementation; (3) rerun every baseline or state for each number where it was taken from; (4) include at least one test set from a source that was not used for training, and do not average across annotation granularities; (5) include a self-configuring baseline and report several seeds; (6) release probability maps, so that any metric can be recomputed. For users: the lesion-level recall for small lesions and the result on an independent source with the user's field of view are more informative than a mean AUPR on IDRiD.

\paragraph{Limitations and scope.} The scope of the study was fixed before the factor experiments were run, to what one protocol and one cluster could carry: the main table has three seeds for nine of its \rn{n.main} configurations and one seed for the others, the factor experiments (training source, granularity, resolution, training amount, loss) use one to five representative models and one seed and report only differences larger than the seed variation of the main table, the encoder comparison has one seed and one learning rate per cell, and the source-recipe track was stopped before any run finished. Bootstrap intervals over test images are given per class (Supplementary Table~4); an interval for the class mean, from the same replicates, is left for the final version, and on both internal tests the per-class intervals of the leading models overlap. Left for future work are the remaining seeds and models of each factor experiment, a registered ranked row for the natural-image encoders under the UPerNet head and the full encoder grid (supervised ImageNet, frozen encoders, learning-rate sweeps, the adaptation-by-initialization square), the second-tier and rewritten models (UNet++, UPerNet, PerSAM, H2Former, L-Seg, SAA, RTNet), the released WSRFNet weights, vessel maps and text as additional inputs, the intermediate settings of each factor (a 1,024-pixel canvas, a quarter of DDR, merging all sources, the bin loss), and the test sets that were not available at the time (e-ophtha, ROC, ultra-widefield OCTA).  Region-consensus and coarse ground truth bias pixel metrics, which is why DiaRetDB1 is reported with its image-level protocol and Retinal-Lesions as a separate column. The licence of FGADR does not allow derived data to be released, and we do not release them for Retinal-Lesions either. Most sources carry no patient identifiers, so leakage can be audited only through near-duplicate images. The ultra-widefield test set comes from a single centre and device. IDRiD has 27 test images, and intervals on it are wide. The first round used the hyperparameters of each source recipe, with a small learning-rate sweep on DDR as the only tuning. H2Former, although registered in tier~1, did not fit the cards used here at its native input size and is not in the main tables. Finally, the frozen lesion-matching rule is permissive (any overlap counts), so lesion-level precision is reported together with a stricter variant in the supplementary tables.
